% =============================================================================
% 5.4 研究前景与展望
% 草稿版本：v0.1 (2026-05-05)
% 字数：约 1100 字
% 风格规则：v0.2 风格（无 \textbf / \textit / 引例括号 / 双引号 / 破折号）
% 引用：\cite{ref1, ref19, ref23, ref29}
% 图表：无
% 依赖宏包：booktabs / tabularx（主稿已加载）
% =============================================================================

\subsection{研究前景与展望}
\label{sec:future-work}

本文的工作是在毕业设计研究范围内对气象 Agent 单点架构、双通路 RAG
机制、代码生成与受限沙箱、报告生成方案进行系统化设计与实证。\ref{sec:scenario-eval}
节的 5 组评测在揭示有效性的同时也暴露了三类局限性，这些局限性自然
指向 4 类未来工作。本节按局限性与未来工作两条线索展开。

\subsubsection*{5.4.1 \quad 局限性回顾}

第一类局限性是多轮长程依赖能力不足，源自 \ref{sec:intent-eval} 节
意图识别评测。\texttt{context\_\bk{}inference} 类 3 条用例端到端通过率仅
为 33.3\%，与单轮场景的 100\% 之间形成 66.7 个百分点差距。当查询
同时切换意图与隐式继承多个槽位时，大语言模型倾向于把所有未在 query
中显式出现的字段都视为重置而非继承。这一现象表明单纯把上下文摘要
拼接到 prompt 不足以实现稳定的多轮槽位继承。

第二类局限性是数值阈值的精确归属仍依赖确定性分级器兜底，源自
\ref{sec:rag-eval} 节通路 A 检索评测。即使在最优混合权重 $\alpha=0.9$
下，60 条评测集上仍有约 9 条 Top-1 脱靶，且全部集中在邻近阈值边界、
术语别名冲突、跨类别多文档三类场景。混合检索的工程上限在 85\% 左右，
难以单纯依靠权重调整突破。本文通过通路 B 即确定性 if-else 分级器
加 \texttt{grade\_\bk{}id} 硬链接的方式覆盖了这一盲区，但本质上是把问题
绕过而非解决。如果未来要扩展到行业气象等复杂领域，分级器规则的
人工编写成本会显著上升。

第三类局限性是受限沙箱的安全性是轻量原型而非生产级方案。
\ref{sec:code-eval} 节的沙箱仅做了 builtins 与模块的白名单约束加
墙钟超时加结构化错误捕获四项，能够阻止意外死循环与文件读写，但不
能抵御主动恶意攻击。生产环境下应当使用 RestrictedPython、Pyodide
或 Docker 容器实现进程级隔离 \cite{ref23}。本文限定在毕业设计研究
范围内主动放弃了这一安全维度，保留为未来工作。

此外，\ref{sec:e2e-eval} 节端到端评测中暴露的集成衰减效应同样值得
列入局限性清单：通路 B 单测 \texttt{citation} 真实性已达 100\%，
但端到端集成后衰减到 26.7\%，意味着大语言模型在最终报告生成阶段
仍可能省略检索到的 GB/T 编号。这是一类只有在集成层评测中才会暴露
的问题，单独优化任何一个模块都无法消除。

\subsubsection*{5.4.2 \quad 未来工作}

针对上述局限性，本文规划 4 类具体的未来工作方向。

第一类是显式对话状态跟踪机制。在 \texttt{recognizer} 与
\texttt{completer} 之外引入独立的 \texttt{DialogState} 数据结构，
每轮对话结束后把当前 \texttt{locations}、\texttt{time}、\texttt{intent}
三个槽位写入状态对象，下一轮 prompt 中以结构化 JSON 形式注入。要求
大语言模型对每个槽位显式输出继承或覆盖的二选一决策，把当前的隐式
推理转化为显式标注。同时把多轮评测集从目前的 3 条扩展到 30 条以上，
覆盖意图切换、地点切换、时间切换、复合切换等典型对话状态切换案例
\cite{ref29}。

第二类是知识库扩展到行业气象领域。本文当前知识库覆盖 GB/T 28592、
GB/T 28591、GB/T 27964、GB 50736、GB/T 21987 等 17 项标准共 62 条
条目，主要面向公众气象服务场景。未来将扩展到航空气象 ICAO Annex 3、
农业气象作物物候期分级、海洋气象浪高等级等行业标准，对应同步扩展
通路 B 分级器规则与通路 A 评测集。同时引入更细粒度的元数据，例如
\texttt{applicable\_\bk{}region} 区域适用性、\texttt{publish\_\bk{}year}
发布年份等，支持随时间漂移的标准更新。

第三类是受限沙箱升级为进程级隔离。具体方案有两条候选路径。一是
RestrictedPython 加 PyPy 沙箱，可以在不引入容器开销的前提下实现
更严格的字节码级访问控制。二是 Pyodide WebAssembly 运行时，把代码
执行隔离到 WASM 沙箱中，天然不暴露任何系统调用接口。两条路径在
延迟、安全性、可移植性三个维度上各有取舍，需要进一步实测对比。

第四类是引用强制注入与生成后校验。针对 \ref{sec:e2e-eval} 节暴露
的集成衰减效应，未来工作包含两个具体改进。一是在 RAG 工具返回结果
后立即把 \texttt{must\_\bk{}cite} 关键字以结构化 JSON 形式注入到模型上下文
中，并在 system prompt 中强制要求最终报告包含该关键字。二是在报告
生成完成后引入独立的引用校验后处理器，对照 \texttt{must\_\bk{}cite}
进行字符串级匹配，未命中时触发模型二次生成。这两类机制的有效性
需要在扩展评测集上独立验证。

进一步的扩展方向包括把单智能体扩展为多智能体协作架构 \cite{ref1, ref19}，
将检索专家、桥接专家、决策专家、报告专家解耦为独立 agent，引入
agent 间共享的 \texttt{Blackboard} 状态空间。多智能体方案在分工
专业化与可调试性上具有优势，但需要解决状态同步与额外延迟两类工程
问题。这一方向超出本文研究范围，保留为后续工作的开放课题。

